% Quarantine-only candidate proof.
% Date: 29 August 2026.
% This file is not a live-book input and has not been Rethlas-verified.
% It makes no novelty or priority claim.

\paragraph{Source conventions.}
We use the type decomposition, branch space, residual curve, and canonical
graph metric of \cite[§§3.3--4.5]{Ducros24CurveStructure}.  In
particular, \cite[§3.3.5.2]{Ducros24CurveStructure} associates to a
type-$2$ point $y$ a projective normal integral residual curve $C_y$ over
the degree-one residue field of $K$; it need not be geometrically
integral.  The branch theorem
\cite[Theorem~4.2.10 and §§4.2.11.1--4.2.11.2]
{Ducros24CurveStructure} applies to an arbitrary analytic curve.  Since
$y$ is interior, its branches are all closed points of $C_y$ in type~$2$,
and there are exactly two branches in type~$3$.

For an analytic unit $f$ on an arbitrary analytic curve, its variation
graph is closed and locally finite, its norm is locally constant off that
graph, and its restriction to the graph is piecewise a dilation by
\cite[§§4.5.21.1--4.5.21.4]{Ducros24CurveStructure}.  The normalized
branch slope $\mathrm d\log_b|f|$ is an integer, without a generic
quasi-smoothness hypothesis.  If the curve is boundaryless, the finite
sum of these slopes at every node is zero by
\cite[Theorem~4.5.22]{Ducros24CurveStructure}.  We use CLD tropical
coordinates with the sign convention $-\log|f|$.

\begin{proof}[Proof of
\cref{lem:p213-variation-graph-branch-separation}]
We first treat the branch separation and then add the proper source
localization.

Suppose that $y$ is of type~$2$, and let $C_y$ be its residual curve.
Every branch $b$ at $y$ corresponds to a unique closed point $P_b$ of
$C_y$.  This assertion needs neither normality nor quasi-smoothness of
$Y$: it is exactly
\cite[Theorem~4.2.10 and §4.2.11.1]{Ducros24CurveStructure}.  Normality
will be used only to lift rational functions on $C_y$ to analytic unit
germs.

Let $A_0\subseteq A$ be the set of branches on which $F_0$ is not
locally constant, and write
\[
                     A_0=\{b_1,\ldots,b_s\},
             \qquad P_i=P_{b_i}.
\]
Weak approximation on the projective normal integral curve $C_y$
provides $r_i\in\operatorname{Rat}(C_y)^\times$ such that
\begin{equation}
       \operatorname{ord}_{P_j}(r_i)=\delta_{ij}
       \qquad(1\leq i,j\leq s).
 \label{eq:p213-variation-first-order-block}
\end{equation}
Here is a Riemann--Roch construction that records the field hypotheses.
Choose an effective divisor $D$, disjoint from the $P_j$, sufficiently
large that
\[
 H^1\!\left(C_y,
   \mathcal O_{C_y}\!\left(D-2P_i-\sum_{j\ne i}P_j\right)\right)=0.
\]
The associated jet-evaluation map is surjective.  Prescribing a
uniformizer modulo $\mathfrak m_{P_i}^2$ at $P_i$ and the value $1$
modulo $\mathfrak m_{P_j}$ at every $P_j\ne P_i$ yields $r_i$.
This uses only that $C_y$ is projective, normal, and integral over its
actual constant field.

Let $B$ be the finite set of closed points outside
$\{P_1,\ldots,P_s\}$ at which at least one $r_i$ has nonzero order.
Choose a fresh effective divisor $D'$, disjoint from
$B\cup\{P_1,\ldots,P_s\}$, sufficiently positive that the simultaneous
first-jet evaluation map at $B$ and the value-evaluation map at the $P_i$
are surjective.  They give $h\in\operatorname{Rat}(C_y)^\times$ such that
\begin{equation}
 \operatorname{ord}_{P_i}(h)=0,
 \qquad
 \operatorname{ord}_{Q}(h)=1
 \quad(Q\in B).
 \label{eq:p213-variation-collision-killer}
\end{equation}
If $R$ is any further zero or pole of $h$, then
$\operatorname{ord}_R(r_i)=0$ for every $i$, by the definition of $B$.
Consequently the order vector
\begin{equation}
 \bigl(\operatorname{ord}_P(r_1),\ldots,
       \operatorname{ord}_P(r_s),\operatorname{ord}_P(h)\bigr)
 \label{eq:p213-variation-order-vector}
\end{equation}
at $P=P_i$ is $(e_i,0)$; at a point of $B$ its last coordinate is one;
and at every further zero or pole of $h$ its first $s$ coordinates are
zero and its last coordinate is nonzero.  Thus the selected rays are
pairwise distinct and no newly activated branch can collide with one of
them.  This single extra function avoids an infinite separator iteration.

We now lift these rational functions.  At a type-$2$ point the field
$\operatorname{Rat}(C_y)$ is the degree-one residue field of
$\mathcal H(y)$.  Choose a norm-one representative of any one of the
functions in \eqref{eq:p213-variation-order-vector} in $\mathcal H(y)$.
The field $\kappa(y)$ is dense in $\mathcal H(y)$; this density is also
used explicitly in the proof of \cite[Lemma~12.2.3]{CLD25}.  Approximate
the representative to error strictly less than one.  The approximant in
$\kappa(y)$ has norm one and the same residue.  Since $Y$ is normal and
$y$ is nonrigid,
\begin{equation}
                         \mathcal O_{Y,y}=\kappa(y)
 \label{eq:p213-variation-normal-local-field}
\end{equation}
by \cite[§4.2.18]{Ducros24CurveStructure}.  Every nonzero approximant
is therefore an analytic unit germ.  After shrinking, all these germs are
analytic units on one common neighborhood.  Denote their tuple by $S$.

We record the saturation factor, because it is not legitimate in a
nowhere-quasi-smooth curve to declare an order-one lift to have normalized
slope one.  Let $b$ correspond to $P$ and let $s_P$ be the lift of a
residual uniformizer at $P$.  Put
\begin{equation}
             e_b:=\operatorname{slope}_b(-\log|s_P|)
                 =-\mathrm d\log_b|s_P|\in\mathbf Z_{>0}.
 \label{eq:p213-variation-branch-index}
\end{equation}
Here is the precise link between the refined valuation and the metric
slope.  Let $\nu_b$ denote the $\mathbf Z$-coordinate of the refined value
at $b$, normalized by $\nu_b(s_P)=1$.  If an analytic unit germ $q$ has
$\nu_b(q)=0$, then \cite[§4.2.6]{Ducros24CurveStructure}, applied both to
$q$ and to $q^{-1}$, shows that $|q|$ is constant on a sufficiently short
section of $b$: the two inequalities supplied there force
$|q|=|q(y)|$.  For a norm-one lift $\widehat g$, the quotient
$\widehat g s_P^{-\operatorname{ord}_P(g)}$ has $\nu_b$-value zero.
Consequently
\begin{equation}
 \operatorname{slope}_b(-\log|\widehat g|)
       =e_b\operatorname{ord}_{P}(g)
       \qquad(g\in\operatorname{Rat}(C_y)^\times),
 \label{eq:p213-variation-residual-slope}
\end{equation}
for any norm-one lift $\widehat g$.  The same argument applied to an
arbitrary analytic unit germ says that its slope is $e_b\nu_b(f)$ and is
therefore divisible by $e_b$.
Since the full moment contains an order-one coordinate, the greatest
common divisor of its slope vector along $b$ is exactly $e_b$.  Hence,
after saturation, its primitive target tangent has residual-coordinate
block equal to the order vector
\eqref{eq:p213-variation-order-vector}, up to the one fixed sign from the
choice of $-\log$.  This proves the asserted ray separation without any
quasi-smoothness or annular-coordinate hypothesis.

If $y$ is of type~$3$, it has exactly two branches because it is interior,
by \cite[§4.2.11.2]{Ducros24CurveStructure}.  A germ whose value at
$y$ has modulus outside $|K^\times|^{\mathbf Q}$ is nonconstant along
both branches.  The node identity
\cite[Theorem~4.5.22]{Ducros24CurveStructure} makes its two normalized
slopes opposite.  Thus the two target rays are already distinct after
appending this one analytic unit germ.  We again denote the resulting
finite separator block by $S$.

Before introducing the boundary block, choose the preceding common
neighborhood to be an open tree $W_*$ around $y$ such that the full
variation graph of the finite block $(F_0,S)$ has no vertex in
$W_*\setminus\{y\}$.  Thus on every arm of $W_* $ the block is either
constant or integral-affine with one constant slope vector, and every
selected $S$-coordinate is strictly monotone on its named arm.  This is
possible by \cite[Theorem~3.5.1 and
§§4.5.21.1--4.5.21.2]{Ducros24CurveStructure}: the union of the finitely
many variation graphs is locally finite, and all arms outside that union
are constant for the whole block.

Apply
\cref{lem:p213-normal-curve-nonrigid-localization} inside $W_*$.  This is
the arbitrary-valuation normal-curve repair of
\cite[Lemma~12.3.6]{CLD25}; in particular, it uses the exact
connected-normal-affinoid irreducibility argument rather than an
irreducibility citation requiring extra hypotheses.  It gives a compact
analytic neighborhood $V\subset W_*$ and a boundary moment
$F_{\mathrm{bd}}$.  Arbitrary positive real radii are allowed in that
construction, so no nontriviality of the ground-field valuation is used.
Concatenating
$F_{\mathrm{bd}}$, $F_0$, and $S$ preserves boundary avoidance by
projection to the first block.  Since $W_*$ is open in $Y$, relative
interiors give $\partial(V/W_*)=\partial(V/Y)$.  The proof of
\cite[Theorem~16.8.5]{CLD25} then gives
\eqref{eq:p213-variation-local-omega}: $U$ is open and the displayed
tropicalization is proper, since it is the base change of the proper map
$F_{\mathrm{trop}}:V\to\mathbf R^N$.

It remains to justify the target and source tubes without altering $V$.
For each coordinate
$f$ of $F$, let $\Gamma_f$ be its variation graph.  By
\cite[§§4.5.21.1--4.5.21.2]{Ducros24CurveStructure}, the finite
union $\Gamma=\bigcup_f\Gamma_f$ is a closed locally finite graph, every
coordinate of $F_{\mathrm{trop}}$ is locally constant off $\Gamma$, and
it is piecewise integral-affine on $\Gamma$.  After the boundary block has
been constructed, subdivide its finitely many local variation arms and
choose the target segment sufficiently close to $v$ that every active arm
there has a constant slope vector.  Its image is then a finite fan at
$v$.  No source domain is shrunk at this stage.  Projection to the
$S$-block and the construction above show that every selected ray is
distinct from the image of every other source branch.  Indeed, a branch
activated by some $r_i$ belongs to $B$ or has a different order vector, a
branch activated only by $h$ has zero $r_i$-block, and a branch activated
only by $F_{\mathrm{bd}}$ has zero $S$-block.  Appending the boundary block
therefore cannot destroy the distinction.

Take a sufficiently short open subsegment $\rho_b$ of a selected ray away
from $v$.  The $S$-coordinate which has order one at its residual point is
strictly monotone on the whole named arm in $W_*$, while the $S$-projection
of every other active arm lies on a distinct ray and every inactive arm
maps to the vertex.  Hence no point on another arm, and no second point on
the named arm, can have the same $S$-value as a point of $\rho_b$.
After $F_{\mathrm{bd}}$ is known, choose $\rho_b$ before the first vertex
of its locally finite variation graph on the named arm.  A subtree attached
therefore either has all coordinates of $F$ constant or has target image
separated from this shortened segment.  Since $F_{\mathrm{trop}}|_U$ is
proper, its image is closed; the preceding $S$-separation then gives a
sufficiently thin ambient open tube $O_b\Subset\Omega$ whose intersection
with $F_{\mathrm{trop}}(U)$ is exactly $\rho_b$.  Every point of
$F_{\mathrm{trop}}^{-1}(O_b)$ therefore lies in the component representing
$b$;
the only additional pieces are subtrees attached to that component on
which all tropical coordinates are constant.  This proves
\eqref{eq:p213-variation-ambient-tube}.  Notice that a ray has empty
ambient interior when $N>1$, so its exact inverse image need not be open.
Proper pushforward is defined on all of $U$; the ambient tube is used only
to isolate the source contribution to one target edge.  No analytic
retraction and no virtual-annulus presentation is involved.
\end{proof}

\begin{proof}[Proof of
\cref{lem:p213-euclidean-one-dimensional-current}]
The assertion is local on $\Omega$.  On a relatively compact coordinate
ball meeting only finitely many cells, positivity gives order zero by
\cite[Proposition~7.2.9]{CLD25}.  Thus $R$ is represented by a symmetric
positive-semidefinite matrix $M=(\mu_{ij})$ of Radon measures: if
$\alpha=\sum\alpha_{ij}\,\mathrm d'x_i\wedge\mathrm d''x_j$, then, up to
the fixed harmless sign convention for currents,
\[
                         \langle R,\alpha\rangle
                    =\sum_{i,j}\int\alpha_{ij}\,\mathrm d\mu_{ij}.
\]
The equation $\mathrm d'R=0$ says, column by column,
\begin{equation}
                         \sum_i\partial_i\mu_{ij}=0
                 \quad\text{in }\mathcal D'(\Omega).
 \label{eq:p213-euclidean-column-divergence}
\end{equation}

Consider first a tube meeting the relative interior of only one edge.
Choose linear coordinates $(t,z_2,\ldots,z_N)$ in which the edge is
$\{z=0\}$.  In
\eqref{eq:p213-euclidean-column-divergence}, use a scalar test function of
the form $z_a\psi(t,z)$, with arbitrary compactly supported $\psi$, for
each $a\geq2$.  Its tangential derivative vanishes on the edge and its
$a$-th normal derivative there is arbitrary.  It follows that
$\mu_{aj}=0$ for every $a\geq2$ and every $j$.  Symmetry gives
$\mu_{ia}=0$ as well.  The remaining equation is
$\partial_t\mu_{11}=0$.  A Radon measure on an interval with zero
distributional derivative is a constant multiple of Lebesgue measure,
and positivity makes the constant nonnegative.  Returning to the integral
affine parametrization by the primitive tangent gives
$R|_e=b_e[e]$ with $b_e\geq0$.

It remains to inspect a vertex $v$.  Subtract the edge currents just found
on a finite star about $v$.  Since the original matrix measures and the
edge measures agree away from $v$, the remainder is
$A_v\delta_v$ for a positive-semidefinite matrix $A_v$.  The distributional
column divergence of the full current at $v$ has the form
\begin{equation}
 \left(\sum_{e\ni v}b_eu_{e,v}\right)\delta_v
       +\left(\sum_i(A_v)_{ij}\partial_i\delta_v\right)_{j=1}^N,
 \label{eq:p213-euclidean-vertex-distribution}
\end{equation}
up to one common sign.  Smooth test functions can have independently
prescribed value and gradient at $v$.  First prescribing value zero and an
arbitrary gradient in
\eqref{eq:p213-euclidean-vertex-distribution} gives $A_v=0$; then
prescribing an arbitrary value gives
\eqref{eq:p213-euclidean-vertex-balance}.  Local finiteness lets these
calculations be made at every vertex and makes the sum in
\eqref{eq:p213-euclidean-edge-decomposition} locally finite.
\end{proof}

\begin{proof}[Proof of
\cref{prop:p213-normal-curve-arbitrary-valuation-current}]
Let $\alpha$ be a compactly supported ordinary CLD-smooth form of
bidegree $(1,1)$ on $Y$.  By the local ordinary-moment presentation
\cite[Proposition~15.1.1(b)]{CLD25} and the dimension argument in the
proof of \cite[Proposition~15.2.1(a)]{CLD25}, its support is contained in
the locus where $d_K(y)\geq1$, which on a curve is $Y^{[2,3]}$.  Thus every
point relevant to a current test is nonrigid, and
the localization in
\cref{lem:p213-variation-graph-branch-separation} applies even when $K$
is trivially valued.  The local moment description
\cite[Proposition~15.1.1(b)]{CLD25}, finite source localization, and the
smooth cutoffs of \cite[Corollary~15.2.9]{CLD25} reduce the test to a
moment $F$ on an open $U$ as in
\eqref{eq:p213-variation-local-omega}.  A cutoff is factored through its
own finite analytic moment by \cite[Corollary~15.1.2]{CLD25}; the cutoff
itself is never treated as an analytic coordinate.  In applying
\cref{lem:p213-variation-graph-branch-separation}, take $A$ to contain
every branch on which the form and cutoff moment blocks vary.  This set is
finite because it is the set of arms at $y$ in a finite union of locally
finite variation graphs.  Additional branches activated by the boundary
block cannot collide with a selected ray, since projection to the
separator block already distinguishes it.

Proper tropical pushforward defines a positive closed Euclidean
supercurrent $T_\Omega$ of bidimension $(1,1)$ supported on the
one-dimensional polyhedral set $F_{\mathrm{trop}}(U)$.  Applying
\cref{lem:p213-euclidean-one-dimensional-current} gives, on the relative
interior of every edge,
\begin{equation}
                         T_\Omega|_e=b_e[e],
             \qquad b_e\geq0,
 \label{eq:p213-variation-edge-density}
\end{equation}
with $b_e$ constant.  At every vertex one has
\begin{equation}
                         \sum_{e\ni v}b_eu_{e,v}=0,
 \label{eq:p213-variation-current-balance}
\end{equation}
where $u_{e,v}$ is the primitive outgoing tangent, and there is no
positive remainder supported at the vertex.

Use the ambient tubes from
\cref{lem:p213-variation-graph-branch-separation} to assign an edge
coefficient to an individual source branch.  Let $m_b>0$ be the
coefficient obtained by pushing forward the canonical integration current
$\delta_Y$, which is closed by \cite[§7.7.4]{CLD25}, and define
\begin{equation}
                             a_b:=b_b/m_b.
 \label{eq:p213-variation-branch-ratio}
\end{equation}
This ratio is independent of the separating moment.  Pass to a common
concatenated moment, and suppose a target projection sends its primitive
source tangent $u'$ to $q u$, where $u$ is primitive and
$q\in\mathbf Z_{>0}$.  Direct parametrization of a calibrated segment
gives
\begin{equation}
                       A_*[u']=q[u].
 \label{eq:p213-variation-segment-pushforward}
\end{equation}
Indeed, the pullback evaluates a $(1,1)$-form on $(qu,qu)$ and the target
parameter is $q$ times the source parameter.  Functoriality of current
pushforward \cite[§7.1.11]{CLD25} therefore multiplies both the coefficient
of $T$ and that of the intrinsic integration current by the same $q$.
The factors cancel in \eqref{eq:p213-variation-branch-ratio}.  The
canonical calibration used for $\delta_Y$ is independent of the moment by
\cite[§§14.3.1--14.3.5]{CLD25}.

At a type-$3$ point there are exactly two branches.  The slope vector of
every analytic unit on one branch is the negative of its slope vector on
the other, by \cite[Theorem~4.5.22]{Ducros24CurveStructure}.  Applying
\eqref{eq:p213-variation-current-balance} first to $T$ and then to
$\delta_Y$ shows that the two branch ratios agree.

Let now $y$ be of type~$2$ and put $C=C_y$.  For every closed point
$P\in C$, let $b_P$ be its branch.  Append a lift of a residual
uniformizer at $P$.  By
\eqref{eq:p213-variation-residual-slope}, saturation divides the full
slope vector by $e_{b_P}$ and makes the residual coordinate equal to one.
The common-refinement argument above then shows that the corresponding
canonical flux $m_P$ is independent of the normalized moment.  Indeed, if
two moments each contain an order-one residual coordinate, their
concatenation still has greatest common slope divisor $e_{b_P}$.  Its
primitive tangent projects to either old primitive tangent with expansion
index one, because the retained order-one coordinate has primitive
coordinate one on both sides.  Thus neither canonical flux changes.
Extend
the two assignments on prime divisors additively:
\begin{equation}
 \begin{split}
 L_T([P])&=a_{b_P}m_P,\\
 L_Y([P])&=m_P.
 \end{split}
 \label{eq:p213-variation-divisor-functionals}
\end{equation}
Both are nonnegative on effective divisors.

Let $g\in\operatorname{Rat}(C)^\times$.  Append a norm-one analytic-unit
lift of $g$, and append the finite separator block from the proof of
\cref{lem:p213-variation-graph-branch-separation} for the branches in
$\operatorname{Supp}(\operatorname{div}(g))$.  The residual-coordinate
component of the primitive tangent in the $P$-branch is
$\operatorname{ord}_P(g)$ by
\eqref{eq:p213-variation-residual-slope}; all other branches have zero
component.  Taking this scalar coordinate in
\eqref{eq:p213-variation-current-balance} gives
\begin{equation}
             L_T(\operatorname{div}(g))=0,
       \qquad L_Y(\operatorname{div}(g))=0.
 \label{eq:p213-variation-principal-zero}
\end{equation}
The second equality is the same calculation for the closed integration
current.  Only finitely many branches enter.

We use Riemann--Roch over the actual base field of $C$; neither geometric
integrality nor a rational point is needed.  Let
$L:\operatorname{Div}(C)\to\mathbf R$ be additive, nonnegative on
effective divisors, and zero on principal divisors.  If $D$ has degree
zero, choose an effective divisor $A$ with
\begin{equation}
                  \deg A+\chi(C,\mathcal O_C)>0.
 \label{eq:p213-variation-rr-positive-constant}
\end{equation}
This is the correct constant over a non-algebraically-closed field:
$H^0(C,\mathcal O_C)$ can be a nontrivial finite extension of the base
field, so $\chi(C,\mathcal O_C)$ must not be replaced by
$1-h^1(C,\mathcal O_C)$.  For every $n\in\mathbf Z$, Riemann--Roch gives
\[
 h^0\bigl(C,\mathcal O_C(A+nD)\bigr)
    \geq \deg A+\chi(C,\mathcal O_C)>0.
\]
Thus there is a nonzero section of
$\mathcal O_C(A+nD)$.  Hence $A+nD$ is linearly equivalent to an
effective divisor, and
\[
                         0\leq L(A)+nL(D)
                         \qquad(n\in\mathbf Z).
\]
Letting $n$ tend to both signs of infinity gives $L(D)=0$.  Thus $L$
factors through the degree subgroup
$\deg\operatorname{Div}(C)\subseteq\mathbf Z$, so
\begin{equation}
                              L=c_L\deg
 \label{eq:p213-variation-degree-factor}
\end{equation}
for one $c_L\geq0$.  Applying this to the two functionals in
\eqref{eq:p213-variation-divisor-functionals} yields
\[
                 m_P=c_Y\deg(P),
       \qquad a_{b_P}m_P=c_T\deg(P).
\]
Since $m_P>0$, one has $c_Y>0$, and all branch ratios at $y$ are equal to
$c_T/c_Y$.  This argument also covers an infinite-valence type-$2$ star,
because every analytic moment and every principal divisor activates only
finitely many branches.

It remains to globalize the constant.  By
\cite[Theorem~3.5.1 and §1.3.5]{Ducros24CurveStructure}, $Y$ is a connected
arcwise-connected graph.  It is not a one-point graph because it is purely
one-dimensional.  Hence it contains an open interval; every interior point
of such an interval belongs to $Y^{[2,3]}$ by normality and
\cite[§4.5.11]{Ducros24CurveStructure}.  Thus this locus is nonempty.  The
same section says that it is convex, so it is path-connected.  Edgewise
constancy from \cref{lem:p213-euclidean-one-dimensional-current}, together
with the type-$2$ and type-$3$ vertex comparisons, propagates one constant
$c$ throughout it.  The localized Euclidean calculation, including the
absence of vertex terms, then gives
$T=c\delta_Y$ on every compactly supported $(1,1)$-test form.  There is no
unseen endpoint term: by the ordinary local presentation and the dimension
argument cited at the start of the proof, every such test form has zero germ
at every point outside $Y^{[2,3]}$.  Branches activated only by the
auxiliary boundary or separator blocks also make no contribution: the
original form moment has zero tangent there, so the pulled-back
$(1,1)$-form vanishes on those edges.  This proves
\eqref{eq:p213-normal-curve-arbitrary-classification}.

Finally, \cite[Example~8.6.4]{CLD25} states that the integration current
is Bedford--Taylor, and the defining identity is stable under
nonnegative scalar multiplication.  Therefore so is $T$.
\end{proof}

\begin{proof}[Proof of
\cref{prop:p213-curve-nonsmooth-locus-dichotomy}]
The quasi-smooth locus is Zariski open by
\cite[§3.1.11]{Ducros24CurveStructure}.  If it is nonempty, its
complement is a proper Zariski closed subset of the irreducible curve and
therefore has dimension zero.  On every strict affinoid domain this
closed analytic subspace is finite: its affinoid algebra is noetherian
and zero-dimensional, hence Artinian.  The strict affinoid
Nullstellensatz, in the form of
\cite[Theorem~4.5]{Tate71Rigid}, makes its points rigid.  These finite descriptions show
that the complement is locally finite in $Y$.  At every interior point,
quasi-smoothness is equivalent to smoothness by
\cite[§3.1.11.4]{Ducros24CurveStructure}; hence it is equivalent to
rig-smoothness when $Y$ is boundaryless.  If the quasi-smooth locus is
empty, the second alternative holds.  This proves the dichotomy and also
shows that a nonrigid failure cannot occur as an isolated exceptional
point.
\end{proof}

\begin{proof}[Verification of
\cref{ex:p213-nowhere-smooth-normal-curve}]
Write $\alpha=t^{1/p}$, so
$L=K[\alpha]/(\alpha^p-t)$.  The algebraic $K$-curve
$\operatorname{Spec}L[T]$ is integral and regular, and its analytification
is good, separated, strict, and boundaryless.  It is normal and purely
one-dimensional.  In characteristic $p$, differentiating the relation
$\alpha^p-t=0$ gives no relation on $\mathrm d\alpha$, and therefore
\begin{equation}
 \Omega^1_{L[T]/K}
      \simeq L[T]\,\mathrm d\alpha
           \oplus L[T]\,\mathrm dT.
 \label{eq:p213-nowhere-smooth-differentials}
\end{equation}
At every analytic point the fibre of the differential sheaf has dimension
two over the completed residue field, whereas the local analytic
dimension is one.  The criterion in
\cite[§3.1.11]{Ducros24CurveStructure} therefore shows that the curve
is nowhere quasi-smooth over $K$.  It contains nonrigid Gauss points, so
its nonsmooth locus is not a locally finite rigid set.
\end{proof}
